% SYNC-ID: threats-to-validity
\section{有效性威胁}\label{sec:threats-to-validity}
% CONTENT_DEFERRED_TO_WRITING_PHASE
\textit{(THIS IS PLACEHOLDER TEXT, PLZ FULFILL IT WHEN WRITING.)}
% SYNC-ID: threats-placeholder-1
有效性限制将在写作阶段根据核验材料记录。本版式占位不声称任何威胁或缓解措施。
% SYNC-ID: threats-placeholder-2
每项未来限制都要配对证据指针和明确的适用范围。
% SYNC-ID: threats-placeholder-3
如果批准的计划需要，可在此处设置内部与外部有效性小节。
% SYNC-ID: threats-placeholder-4
缓解措辞须与复核时可用的证据强度相称。
% SYNC-ID: threats-placeholder-5
本段预留复现与可迁移性说明，但不做任何预测。
% SYNC-ID: threats-placeholder-6
作者核对措辞和来源后，才会列出未解决风险。
% SYNC-ID: threats-placeholder-7
本节预留一句收束语，把每项记录的限制连接到它所限定的解释。只有相应证据指针可用时才会补入正式文字。
% SYNC-ID: threats-placeholder-8
最后一处版式位置用于区分已有记录的限制和仍待写作的任务。该区分将依据与英文伴随文本相同的来源记录核对。
% SYNC-ID: threats-placeholder-8b
有效性章节后续可在此说明某项限制影响解释、迁移性还是可复现性。正式句子必须标明边界，不能暗示未经检查的补救措施。当前文字刻意保持中性，只用于让双栏页面的密度自然。
